% !TeX root = ../PIFB2.tex
\section{Physical Targets: Gravity, Persistence and Measurement}\label{sec:speculative_physics}

\subsection{What would turn a geometric readout into gravity?}
\label{sec:gravity_pullback}

A curved statistical tensor is not yet a gravitational field. The covariant Fisher pullback is positive semidefinite; the separate Lorentz construction supplies an indefinite tensor only after a complex channel, soldering and scale have been chosen. To interpret a resulting Lorentz metric $\gamma$ gravitationally, the same collective construction must also specify the matter or record variables that couple to it and an evolution law for that coupling. Otherwise one can fit a geometry to observations without explaining how geometry affects the participating systems.

A conditional route is available in the normalized population model. Suppose a proposed collective geometry changes a prior, a record kernel or an effective action. Its dependence then enters the evidence and all variational derivatives through that normalized law. For example, if a conditional density is written as $p_\gamma(y)=e^{-S_\gamma(y)}/Z_\gamma$ against a fixed reference measure and differentiation under the integral is justified, then
\begin{equation}
\delta\log Z_\gamma=-\mathbb E_{p_\gamma}\delta S_\gamma,
\qquad
\delta\log p_\gamma=-\delta S_\gamma+\mathbb E_{p_\gamma}\delta S_\gamma.
\label{eq:geometry_feedback_normalizer}
\end{equation}
If the measure itself depends on $\gamma$, its logarithmic variation belongs in the same calculation. A proposed geometric equation that varies only $S_\gamma$ and discards this term is not the equation of the declared probability model.

The physical target is a common causal geometry with consistent source coupling, conservation identities, and an observational limit. In a weak and slowly varying regime, a candidate must predict the appropriate acceleration and clock comparisons from its retained dynamics; naming a curvature scalar or a precision parameter does not supply those predictions. No derivation of Einstein dynamics or its constants is asserted here. The exact coarse effective action of Section~\ref{sec:rg_status} provides a place to investigate induced interactions and their symmetries, while the Lorentz construction specifies the geometric type such a limit would have to support.

\subsection{Persistence and the meaning of an object}

The formation theory distinguishes coherence, persistence and autonomy. The proposed physical reading should keep the distinction. A cluster can agree on a communicated statistic while its private state changes; a long-lived memory can retain that statistic while failing to support a second formation episode; and an exact instantaneous marginal can fail to have closed future dynamics. Calling the cluster a macroscopic object requires evidence for its retained state and response across those changes, rather than only a small divergence at one time.

Precision by itself does not solve this problem. In the scalar composite consensus specialization, take a receiver $q_i=\mathcal N(a,v)$ with fixed $v>0$ and a source $q_j=\mathcal N(0,\delta)$, $\delta>0$. Their cost is
\begin{equation}
K_{ij}=\tfrac12\left[\frac{v+a^2}{\delta}-1+\log\frac{\delta}{v}\right].
\label{eq:precision_source_cost}
\end{equation}
At fixed nonzero coupling, the mean gradient $\partial_aK_{ij}=a/\delta$ becomes strong as $\delta$ decreases. With Gibbs source selection, however, a competing source of bounded cost gives
\begin{equation}
\beta_{ij}\leq C e^{-K_{ij}/\tau},\qquad
\lim_{\delta\downarrow0}\beta_{ij}=0
\label{eq:precision_attention_rejection}
\end{equation}
for fixed $\tau>0$ and fixed positive prior weights. The bound follows by retaining the competing source's positive term in the softmax denominator. Its exponential suppresses the $1/\delta$ growth of the gradient as well. A highly precise but incompatible source can therefore be ignored instead of becoming a universal anchor. A noisy interface or an adaptive receiver changes the calculation and must be modeled explicitly.

Appendix~\ref{app:examples} develops the response and memory calculations relevant to inertia and persistence. Their physical use requires a retained observable, an external perturbation, and an independently identified response law. Curvature of a free-energy landscape supplies stiffness near a stationary state; a kinetic coefficient, dissipation and a memory kernel have different roles. Assigning all of them the name mass would obscure the mechanism. The useful question is whether eliminating internal degrees of freedom produces a reproducible effective response that remains valid over the claimed time window and through another coarse-graining step.

\subsection{Measurement as a recorded interaction}
\label{sec:measurement_analogy}

In the normalized theory a measurement is an interaction with a specified record kernel. The record may be endogenous to the agent network. Its status as a recorded value is distinct from the agents' later beliefs about it, and several agents reading the same record does not multiply the likelihood by the number of readers. This distinction prevents agreement from being manufactured by counting one observation repeatedly.

A classical account of public measurement can seek three concrete properties: the record is stable on the readout timescale; different access channels recover compatible predictions from it; and interventions distinguish the inferred system from alternative explanations. The source of stability could be a metastable collective law or retained memory, but it must be established in the chosen dynamics. Merely increasing a Gaussian precision, or obtaining identical posterior means, proves none of these properties in general.

The present probability model has no noncommuting measurement algebra, Born rule or quantum instrument. Hermitian classical moments in Section~\ref{sec:signature_resolution} do not add those structures. Consequently neither consensus nor the Lorentz representation resolves quantum measurement. A quantum extension would need to specify its state space, composition rule and operational probabilities before asking whether its classical, coarse-grained records reproduce the present theory. This preserves a physical research target while avoiding a claim of quantum emergence from notation alone.

